% Answers to Chapter 4's exercises, lectures/L04/appendix/b_exercises.md. The book's own answers:
% the lecture publishes none.
\answerchapter{4}{Analog-to-Digital Conversion}

\answer{c4:ex:1}
Invert the chapter's mapping, \code{result = 1023 * (V_in / V_ref)}:
\[
  V_\text{in} = \frac{\text{reading} \times V_\text{ref}}{1023}
    = \frac{768 \times 5\,\text{V}}{1023} = \frac{3840}{1023}\,\text{V} \approx 3.754\,\text{V}.
\]
Check: $1023 \times 3.754 / 5 = 768.1$, which rounds back to 768. The datasheet's form, with 1024,
gives $768 \times 5 / 1024 = 3.750$\,V, less than one step (about 5\,mV) away.

\answer{c4:ex:2}
The ADC is a successive-approximation converter. It samples the input onto a small capacitor, then
decides the result one bit per ADC clock cycle, comparing the held voltage against a trial voltage
generated inside the converter, and its analog circuitry needs time to settle at each step. The
datasheet therefore asks for an ADC clock between 50 and 200\,kHz for full 10-bit resolution. The
CPU runs at 16\,MHz, 80 times the upper limit, and slowing the whole CPU down to suit one
peripheral is not an option; so the ADC has a prescaler of its own, the \code{ADPS2:0} bits of
\code{ADCSRA}, which divides the system clock for the ADC alone. At 16\,MHz, \textdiv128
(125\,kHz) is the only setting that lands in range; \textdiv64 already gives 250\,kHz.

At the full 16\,MHz a conversion would finish sooner, but accuracy would collapse: the sampling
capacitor and the comparator would get far too little time to settle, so the lower bits would turn
into noise, readings would jump around and drift from the true voltage, and at 80 times the limit
little of the 10 bits would be left. (The prescaler cannot in fact divide by less than 2, so
8\,MHz is the fastest the ADC can be clocked, still 40 times too fast.)

\answer{c4:ex:3}
\code{MUX0}, bit 0 of \code{ADMUX}. The channel number is written into \code{MUX3:0}, the low four
bits: \code{A0} is channel 0 (\code{0000}) and \code{A1} is channel 1 (\code{0001}), so moving from
\code{A0} to \code{A1} sets \code{MUX0} and leaves \code{REFS1:REFS0}, the \code{AVCC} reference,
as they are. With \code{adc_read()} that is simply a call with \code{channel} set to 1.

\answer{c4:ex:4}
Setting \code{ADSC} only \emph{starts} a conversion. The conversion takes 13 ADC clock cycles, or
25 for the first one after the ADC is enabled; at 125\,kHz that is 104\,\textmu s or 200\,\textmu s,
1,664 or 3,200 CPU cycles, while the CPU reaches the next statement after one or two. Until the
conversion ends, \code{ADC} still holds the previous conversion's result, or 0 if there has been
none, so reading it straight away returns an old value, not the one just asked for. The hardware
sets \code{ADIF} at the moment the conversion completes and its result is in
\code{ADCL}/\code{ADCH}, so waiting for \code{ADIF} is waiting for the result that belongs to this
conversion. (\code{ADSC} reads as 1 while a conversion is in progress, so waiting for it to clear
works too, and the ADC interrupt does the same without polling.)

\answer{c4:ex:6}
It would not crash; it would quietly return 0 from every read. The reason is the statement that
follows the wait: \code{ADCSRA = (1U << ADIF);} writes the whole register, not just the flag, and
every bit written as 0 is cleared, \code{ADEN} included. Writing \code{ADEN} to zero switches the
ADC off, and the datasheet is explicit that switching it off while a conversion is in progress
terminates that conversion. Without the wait, this statement runs a few CPU cycles after the
conversion was started, about 3,200 cycles before it could have finished (a conversion started
together with \code{ADEN} takes 25 ADC clock cycles). So every conversion is aborted: none ever
completes, sets \code{ADIF}, or writes \code{ADCL}/\code{ADCH}. \code{return ADC;} then reads a data
register that no conversion has ever written, and returns its reset value, 0, on every call,
whatever the potentiometer says.

Nothing illegal happens and the program runs exactly as before, only on wrong data. In the program
of \secref{c4:a:4}, every duty cycle is 0, \code{on_ms} is 0 and \code{off_ms} is 10, so the LED
stays dark whenever reading is enabled, which looks more like a wiring fault than a software one.
(Had the flag been cleared without switching the ADC off, the wait's absence would show up
differently: each read would return the result of the conversion before it, one reading behind,
which is stale data that is even harder to spot.)
